\subsection{The $L^1$ seminorm and approximation}

The content of this subsection is also known, and corresponds to \S4 and \S6 of
\cite{BM}.
We give proofs because it is used essentially in
Section~\ref{chap:operator}.

\begin{definition}\label{def:l1-seminorm}
  Let $(X,L,I)$ be an integration system on a set $X$.
  For each $f\in L$ we put
  \[
    \|f\|_I:=I(|f|) .
  \]
  Since $L$ is a Riesz space we have $|f|\in L$, so this value is defined.
\end{definition}

\begin{lemma}\label{lem:l1-seminorm}
  $\|\cdot\|_I$ is a seminorm on $L$. That is, for all $f,g\in L$ and
  $a\in\mathbb{R}$,
  \[
    \|af\|_I=|a|\,\|f\|_I,
    \qquad
    \|f+g\|_I\le\|f\|_I+\|g\|_I.
  \]
\end{lemma}
\begin{proof}
  The first identity follows from $|af|=|a||f|$ and linearity of $I$.
  Moreover $|f+g|\le|f|+|g|$ and
  Lemma~\ref{lem:integration-preserve-order} give
  $I(|f+g|)\le I(|f|)+I(|g|)$.
\end{proof}

\begin{remark}
  In general $\|f\|_I=0$ does not imply $f=0$.
  The operation that collapses this discrepancy is treated in
  Subsection~\ref{sec:ae}.
\end{remark}

\begin{lemma}[approximation of integrable functions by $L$]\label{lem:l1-approx-by-source}
  Let $(X,L,I)$ be an integration system on a set $X$ and let
  $(X,L^1,I^1)$ be its Daniell extension.
  Then for every $f\in L^1$ there is a sequence of functions
  $\{s_N\}\subset L$ satisfying
  \[
    \|f-s_N\|_{I^1}=I^1(|f-s_N|)\longrightarrow0
    \qquad(N\to\infty) .
  \]
  That is, $L$ is dense in $L^1$ with respect to $\|\cdot\|_{I^1}$.
\end{lemma}
\begin{proof}
  Choose a series representation of $f$, say
  \[
    f\overset{\varphi_n}{\simeq}\sum_{n=1}^{\infty}f_n ,
  \]
  and put
  \[
    s_N:=\sum_{n=1}^Nf_n\in L,
    \qquad
    r_N:=f-s_N .
  \]
  Since finitely many values $\varphi_n(x)$ are always finite,
  \[
    r_N\overset{\varphi_{N+m}}{\simeq}
    \sum_{m=1}^{\infty}f_{N+m}
  \]
  is a series representation of $r_N$.

  For each $m\ge1$ put
  \[
    g_{N,m}:=
    \left|\sum_{k=1}^{m}f_{N+k}\right|
    -\left|\sum_{k=1}^{m-1}f_{N+k}\right| ,
  \]
  where the empty sum is $0$. Since $L$ is a Riesz space we have
  $g_{N,m}\in L$, and the triangle inequality for absolute values gives
  \[
    |g_{N,m}|\le |f_{N+m}|\le\varphi_{N+m} .
  \]
  Moreover, if $\sum_m\varphi_{N+m}(x)<+\infty$ then, by telescoping the sum and
  continuity of the absolute value,
  \[
    \sum_{m=1}^{\infty}g_{N,m}(x)
    =\left|\sum_{m=1}^{\infty}f_{N+m}(x)\right|
    =|r_N(x)| .
  \]
  Hence
  \[
    |r_N|\overset{\varphi_{N+m}}{\simeq}
    \sum_{m=1}^{\infty}g_{N,m}
  \]
  is a series representation of $|r_N|$. By positivity of the integral,
  \[
    \begin{aligned}
      0\le I^1(|f-s_N|)
      &=\sum_{m=1}^{\infty}I(g_{N,m})\\
      &\le\sum_{m=1}^{\infty}I(|g_{N,m}|)\\
      &\le\sum_{m=1}^{\infty}I(\varphi_{N+m})
      \longrightarrow0 ,
    \end{aligned}
  \]
  which is the assertion.
\end{proof}

\begin{corollary}\label{cor:tail-estimate}
  Let $(X,L,I)$ be an integration system and suppose $\{h_j\}\subset L$
  satisfies $\sum_j I(|h_j|)<+\infty$.
  Then the function
  \[
    h(x):=\begin{cases}
      \sum_j h_j(x) & \text{if $\sum_j|h_j(x)|<+\infty$}\\
      0 & \text{otherwise}
    \end{cases}
  \]
  belongs to $L^1$, and for each $J\ge1$
  \[
    \left\|h-\sum_{j=1}^{J}h_j\right\|_{I^1}\le\sum_{j>J}I(|h_j|)
  \]
  holds.
\end{corollary}
\begin{proof}
  Putting $\varphi_j:=|h_j|\in L$, the pair $\{h_j\},\{\varphi_j\}$ is a series
  representation of $h$, so $h\in L^1$.
  Moreover $h-\sum_{j\le J}h_j$ has $\{h_{J+m}\}_{m\ge1}$ and
  $\{\varphi_{J+m}\}_{m\ge1}$ as a series representation, so the same estimate
  as at the end of the proof of Lemma~\ref{lem:l1-approx-by-source} gives the
  stated inequality.
\end{proof}

\begin{theorem}[$\|\cdot\|$ completeness]\label{thm:l1-norm-complete}
  Let $(X,L,I)$ be a Daniell closed integration system.
  Then $(L,\|\cdot\|_I)$ is complete.
  That is, for every Cauchy sequence $\{g_k\}\subset L$ with respect to
  $\|\cdot\|_I$ there is a $g\in L$ with $\|g-g_k\|_I\to0$.
\end{theorem}
\begin{proof}
  Let $\{g_k\}$ be a Cauchy sequence. Choose a subsequence $\{g_{k_j}\}$ so that
  \[
    \|g_{k_{j+1}}-g_{k_j}\|_I<2^{-j} ,
  \]
  and put
  \[
    h_1:=g_{k_1},\qquad h_{j+1}:=g_{k_{j+1}}-g_{k_j} .
  \]
  Then $\sum_j I(|h_j|)<+\infty$.
  By Corollary~\ref{cor:tail-estimate} there is a corresponding $h\in L^1$ with
  \[
    \left\|h-g_{k_J}\right\|_{I^1}
    =\left\|h-\sum_{j=1}^{J}h_j\right\|_{I^1}
    \le\sum_{j>J}I(|h_j|)
    \longrightarrow0 .
  \]
  Since $(X,L,I)$ is Daniell closed we have $L^1=L$, and
  $I^1|_L=I$ by Lemma~\ref{lem:daniell-ext-contains-source}.
  Hence $h\in L$ and $\|h-g_{k_J}\|_I\to0$.
  A Cauchy sequence with a convergent subsequence converges to the same limit as
  a whole, so $\|h-g_k\|_I\to0$.
\end{proof}
